\documentclass[a4paper]{article}

\usepackage[margin=1in]{geometry}
\usepackage{amsmath,amsfonts,mathtools,amsthm,amssymb}
\usepackage[l2tabu,orthodox]{nag}
\usepackage{microtype} % fixes spacing or whatever
\usepackage{enumitem}
\usepackage{tikz-cd}
\usepackage{xcolor}
\usepackage{physics}
\usepackage{mathrsfs}
\usepackage{cancel}
% \usepackage{libertine,libertinust1math}
\usepackage{eucal}
\usepackage[toc,page]{appendix}
\usepackage{pgfplots}
\pgfplotsset{compat=1.18}

\newtheorem{proposition}{Proposition}
\newtheorem{theorem}{Theorem}
\newtheorem{lemma}{Lemma}
\newtheorem{corollary}{Corollary}

\theoremstyle{definition}
\newtheorem{remark}{Remark}
\newtheorem{definition}{Definition}
\newtheorem{example}{Example}

% \usepackage[sc,osf]{mathpazo}   % With old-style figures and real smallcaps.
% \linespread{1.025}              % Palatino leads a little more leading
% % Euler for math and numbers
% \usepackage[euler-digits,small]{eulervm}
% \AtBeginDocument{\renewcommand{\hbar}{\hslash}}

\usepackage{etoolbox}
\usepackage{dynkin-diagrams}
\def\row#1/#2!{\dynkin{#1}{#2}\\}
\newcommand{\tble}[1]{
   \renewcommand*\do[1]{\row##1!}
   \[
      \begin{array}{ll}\docsvlist{#1}\end{array}
   \]
}

\allowdisplaybreaks

\renewcommand{\epsilon}{\varepsilon}
% \renewcommand{\phi}{\varphi}

\DeclareMathAlphabet{\mathpzc}{OT1}{pzc}{m}{it}

\newcommand{\goth}[1]{\mathfrak{#1}}
\newcommand{\red}[1]{#1_{\text{red}}}
\newcommand{\ad}[1]{#1_{\text{Ad}}}
\newcommand{\reg}[1]{#1_{\text{reg}}}
\newcommand{\sh}[1]{#1^{\text{sh}}}
\newcommand{\cpdv}[2]{\cfrac{\partial #1}{\partial #2}}

\newcommand{\fA}{\mathcal{A}}
\newcommand{\fB}{\mathcal{B}}
\newcommand{\fC}{\mathcal{C}}
\newcommand{\fD}{\mathcal{D}}
\newcommand{\fE}{\mathcal{E}}
\newcommand{\fF}{\mathcal{F}}
\newcommand{\fG}{\mathcal{G}}
\newcommand{\fH}{\mathcal{H}}
\newcommand{\fI}{\mathcal{I}}
\newcommand{\fJ}{\mathcal{J}}
\newcommand{\fK}{\mathcal{K}}
\newcommand{\fL}{\mathcal{L}}
\newcommand{\fM}{\mathcal{M}}
\newcommand{\fN}{\mathcal{N}}
\newcommand{\fO}{\mathcal{O}}
\newcommand{\fP}{\mathcal{P}}
\newcommand{\fQ}{\mathcal{Q}}
\newcommand{\fR}{\mathcal{R}}
\newcommand{\fS}{\mathcal{S}}
\newcommand{\fT}{\mathcal{T}}
\newcommand{\fX}{\mathcal{X}}
\newcommand{\fY}{\mathcal{Y}}
\newcommand{\fZ}{\mathcal{Z}}

\newcommand{\be}{\mathbf{e}}
\newcommand{\bx}{\mathbf{x}}
\newcommand{\bone}{\mathbf{1}}
\newcommand{\cx}{\bar{x}}
\newcommand{\cy}{\bar{y}}
\newcommand{\vx}{\vec{x}}
\newcommand{\vsigma}{\vec{\sigma}}
\newcommand{\tzeta}{\tilde{\zeta}}

\newcommand{\gm}{\goth{m}}
\newcommand{\gp}{\goth{p}}
\newcommand{\gF}{\goth{F}}
\newcommand{\gU}{\goth{U}}
\newcommand{\gV}{\goth{V}}

\newcommand{\A}{\mathbf{A}}
\newcommand{\C}{\mathbf{C}}
\newcommand{\F}{\mathbf{F}}
\newcommand{\G}{\mathbf{G}}
\newcommand{\bH}{\mathbf{H}}
\newcommand{\N}{\mathbf{N}}
\newcommand{\PP}{\mathbf{P}}
\newcommand{\Q}{\mathbf{Q}}
\newcommand{\R}{\mathbf{R}}
\newcommand{\Z}{\mathbf{Z}}

\newcommand{\dlog}{\dd\log}

\newcommand\srestr[2]{{\left.\kern-\nulldelimiterspace #1\vphantom{\small|} \right|_{#2}}}
\newcommand\restr[2]{{\left.\kern-\nulldelimiterspace #1 \vphantom{\big|} \right|_{#2}}}
\newcommand\gsec[2]{\Gamma({#1}, {#2})}
\newcommand\gsecx[1]{\Gamma(X, {#1})}

\DeclareMathOperator{\chr}{char}
\DeclareMathOperator{\rProj}{\mathbf{Proj}}
\DeclareMathOperator{\rSpec}{\mathbf{Spec}}
\DeclareMathOperator{\rHom}{\mathbf{Hom}}
\DeclareMathOperator{\rExt}{\mathbf{Ext}}
\DeclareMathOperator{\id}{id}
\DeclareMathOperator{\Frac}{Frac}
\DeclareMathOperator{\rk}{rank}
\DeclareMathOperator{\deter}{det}
\DeclareMathOperator{\pic}{Pic}
\DeclareMathOperator{\cacl}{CaCl}
\DeclareMathOperator{\trd}{tr.d.}
\DeclareMathOperator{\cl}{Cl}
\DeclareMathOperator{\depth}{depth}
\DeclareMathOperator{\codim}{codim}
\DeclareMathOperator{\Div}{Div}
\DeclareMathOperator{\coker}{coker}
\DeclareMathOperator{\len}{length}
\DeclareMathOperator{\height}{height}
\DeclareMathOperator{\supp}{Supp}
\DeclareMathOperator{\proj}{Proj}
\DeclareMathOperator{\im}{im}
\DeclareMathOperator{\Hom}{Hom}
\DeclareMathOperator{\Der}{Der}
\DeclareMathOperator{\spec}{Spec}
\DeclareMathOperator{\Aut}{Aut}
\DeclareMathOperator{\ch}{char}
\DeclareMathOperator{\tor}{Tor}
\DeclareMathOperator{\Ann}{Ann}
\DeclareMathOperator{\Syl}{Syl}
\DeclareMathOperator{\Sym}{Sym}
\DeclareMathOperator{\GL}{GL}
\DeclareMathOperator{\SL}{SL}
\DeclareMathOperator{\Stab}{Stab}
\DeclareMathOperator{\Perm}{Perm}
\DeclareMathOperator{\Orb}{Orb}
\DeclareMathOperator{\Gal}{Gal}
\DeclareMathOperator{\Supp}{Supp}
\DeclareMathOperator{\Ext}{Ext}
\DeclareMathOperator{\Tor}{Tor}
\DeclareMathOperator{\PGL}{PGL}
\DeclareMathOperator{\hd}{hd}
\DeclareMathOperator{\pd}{pd}
\DeclareMathOperator{\SO}{SO}
\DeclareMathOperator{\SU}{SU}
\DeclareMathOperator{\Sp}{Sp}
\DeclareMathOperator{\Oth}{O}
\DeclareMathOperator{\Uni}{U}
\DeclareMathOperator{\evf}{ev}
\DeclareMathOperator{\cH}{H}
\DeclareMathOperator{\dR}{R}
\DeclareMathOperator{\End}{End}
\DeclareMathOperator{\Hasse}{Hasse}
\DeclareMathOperator{\Num}{Num}

\author{James Lee}

% Solarized
% \pagecolor[RGB]{0,20,26}
% \color[RGB]{191,191,191}

% Sephia
% \pagecolor[RGB]{249,239,220}

% Grey
% \pagecolor[RGB]{200, 200, 200}

% Black
% \pagecolor[RGB]{0,0,0}
% \color[RGB]{255,255,255}


\title{Chapter 5, Section 3}

\begin{document}

\maketitle

\begin{enumerate} [label=\textbf{\arabic*.}, leftmargin=0em]

\item Let $X$ be a nonsingular projective variety of any dimension, let $Y$ be a nonsingular subvariety, and let $\pi : \tilde{X} \to X$ be obtained by blowing up $Y$.
Show that $p_a(\tilde{X}) = p_a(X)$.

\begin{proof}
   We can generalize the proof of (3.4). 

   \begin{proposition}
        In the above setting, we have $\pi_*\fO_{\tilde{X}} = \fO_X$, and $\dR^i \pi_*\fO_{\tilde{X}} = 0$ for $i > 0$.
        Therefore, $\cH^i(X, \fO_X) \cong \cH^i(\tilde{X}, \fO_{\tilde{X}})$ for all $i \geq 0$.
   \end{proposition}

   By (II, 8.24), the strict transform of $Y$ in $\tilde{X}$ is isomorphic to the projection space bundle $\P(\fI/\fI^2)$, where $\fI$ is the ideal sheaf $f$.
   It is locally free of rank $r = \codim(Y, X)$ since $Y$ is nonsingular.
   In particular, $\P(\fI/\fI^2)$ has dimension $Y + \codim{r} - 1$, so it corresponds to a divisor on $X$.
   By (II, 7.15), there exists a commutative diagram
   \[ \begin{tikzcd}
\tilde{Y} \arrow[r, "\tilde{f}"] \arrow[d] & \tilde{X} \arrow[d] \\
Y \arrow[r, "f"]                           & X                  
\end{tikzcd} \]
    where $f$ is a closed immersion, and so is $\tilde{f}$.
    By the theorem on formal functionns, For any scheme theoretic point $y \in Y$,  the natural map
    \begin{equation*}
        \hat{\fF^i}_y = (\dR^i \pi_* \fO_{\tilde{X}})_y \otimes \widehat{\fO}_{y, Y} \cong \varprojlim \cH^i(Y_n, \fO_{Y_n}).
    \end{equation*}
    By (7.11), $\pi_*\fO_{\P(\fI/\fI^2)} = \fO_X$

   Since $\pi$ is an isomorphism of $V = \tilde{X} - E$ onto $U = X - Y$ (II, 7.13).
   The sheaves $\fF^i = \dR^i \pi_* \fO_{\tilde{X}}$ for $i > 0$ have support at $Y$.
\end{proof}

\item Let $C$ and $D$ be curves on a surface $X$, meeting at a point $P$.
Let $\pi : \tilde{X} \to X$ be the monoidal transformation with center $P$.
Show that $\tilde{C} . \tilde{D} = C. D - \mu_P(C) \cdot \mu_P(D)$.
Conclude that $C.D. = \sum \mu_P(C) \cdot \mu_P(D)$, where the sum is taken over all intersection points of $C$ and $D$, including infinitely near intersection points.

\item Let $\pi : \tilde{X} \to X$ be a monoidal transformation, and let $D$ be a very ample divisor on $X$.
Show that $2\pi^*D - E$ is ample on $\tilde{X}$.

\item \textit{Multiplicity of a Local Ring.} Let $A$ be a Noetherian local ring with maximal ideal $\gm$.
For any $l > 0$, let $\psi(l) = \len(A/\gm^i)$.
We call $\psi$ the \textit{Hilbert-Samuel function} of $A$.
\begin{itemize}
    \item[(a)] Show that there is a polynomial $P_A(z) \in \Q[z]$ such that $P_A(l) = \psi(l)$ for all $l \ag 0$.
    This is the \textit{Hilbert-Samuel polynomial} of $A$.
    \item[(b)] Show that $\deg{P_A} = \dim{A}$.
    \item[(c)] Let $n = \dim{A}$.
    Then we define the \textit{multiplicity} of $A$, denoted $\mu(A)$, to be $(n!) \cdot $ leading coefficient of $P_A$.
    If $P$ is a point on a Noetherian scheme $X$, we define the \textit{multiplicity} of $P$ on $X$, $\mu_P(X)$, to be $\mu(\fO_{P, X})$.
    \item[(d)] Show that for a point $P$ on a curve $C$ on a surface $X$, this definition of $\mu_P(C)$ coincides with the one in the text just before (3.5.2).
    \item[(e)] If $Y$ is a variety of degree $d$ in $\P^n$, show that the vertex of the cone over $Y$ is a point of multiplicity $d$.
\end{itemize}

\item Let $a_1, \dots, a_r$, $r \geq 5$, be distinct elements of $k$, and let $C$ be the curve in $\P^2$ given by the (affine) equation $y^2 = \prod_{i = 1}^r (x - a_i)$.
Show that the point $P$ at infinity on the $y$-axis is a singular point.
Compute $\delta_P$ and $g(\tilde{Y})$, where $\tilde{Y}$ is the normalization of $Y$.
Show in this way that one obtains hyperelliptic curves of ever genus $g \geq 2$.

\item Show that analytically isomorphic curve singularities (I, 5.6.1) are equivalent in the sense of (3.9.4), but not conversely.

\item For each of the following singularities at $(0, 0)$ in the plane, give an embedded resolution, compute $\delta_P$, and decide which ones are equivalent.
\begin{itemize}
    \item[(a)] $x^3 + y^5 = 0$
    \item[(b)] $x^3 + x^4 + y^5 = 0$
    \item[(c)] $x^3 + y^4 + y^5 = 0$
    \item[(d)] $x^3 + y^5 + y^6 = 0$
    \item[(e)] $x^3 + xy^3 + y^5 = 0$
\end{itemize}

\begin{proof} For each $f(x, y) = 0$, denote $F(x, y, z) = 0$ the projectivization.
    For each homogenous degree $r$ polynomial $F(x, y, z) \in k[x, y, z]$, we compute the singular locus of $F$ as the ideal $J = (F, \partial_x F, \partial_y F, \partial_z F)$, where
    \begin{equation*}
        \partial_x F := \frac{\partial F}{\partial x}, \quad \partial_y F := \frac{\partial F}{\partial y}, \quad \partial_z F := \frac{\partial F}{\partial z}
    \end{equation*}
    Notice that $\mu_P(f) = 3$ for all $f$ above.
    In particular, $\delta_P = 3$.

\begin{itemize} [leftmargin=0cm]
    \item[(a)]
    The curve is parametrized by
    \begin{equation*}
        z \mapsto (z^5, -z^3).
    \end{equation*}
    The tangent vector vanishes at an order of 2, which means $\mu_P(C) = 3$.
    
    Blowing up in coordinates $x = ty$ gives
    \begin{equation*}
        y^3 (t^3 + y^2) = 0,
    \end{equation*}
    where the exceptional divisor $E_1$ is defined by $x = y = 0$.
    Then $\pi^*C = \tilde{C} + 3E$.
    Since $\tilde{C}$ is still singular, parametrized by
    \begin{equation*}
        z \mapsto (x, y, t) = (iz^5, iz^3, z^2).
    \end{equation*}
    The tangent vector vanishes at an order of 1, so we have $\mu_P(\tilde{C}) = 2$, so $\delta_P = 1$.
    It is either the cusp or node, 

    \item[(b)] 
    

    \item[(c)]

    \item[(d)] $t^3+y^2+y^3=0$, cusp

    \item[(e)] $t^3 + ty + y^2 = 0$, node
\end{itemize}
\end{proof}

\item[\textbf{8.}] Show that the following two singularities have the same multiplicity, and the same configuration of infinitely near singular points with the same multiplicities, hence the same $\delta_P$, but are not equivalent.
\begin{itemize}
    \item[(a)] $x^4 - xy^4 = 0$
    \item[(b)] $x^4 - x^2 y^3 - x^2 y^5 + y^8 = 0$
\end{itemize}

\end{enumerate}

\end{document}
